% lmcs, 2018.5.2, 2019.9.9, 2020.2.10
\documentclass{lmcs}
%\usepackage{fouriernc}
%\setmathfont[FakeBold=0.5]{Latin Modern Math}

\usepackage{wasysym}
\usepackage{hyperref}
\usepackage{graphicx}
\usepackage{amssymb,amsmath,graphicx,color}
\usepackage{mymath,proof,latexsym}
\usepackage{xcolor}
\usepackage[pdftex,outline]{contour}
\usepackage{tikz}
%%%%%%%%%%%%%%%%%%%%%%%%%%%%%%%%%%%%
% WE ADD THE COMMANDS 
% theorem lemma proposition corollary definition  
%%%%%%%%%%%%%%%%%%%%%%%%%%%%%%%%%%%%
\newtheorem{theorem}{Theorem}[section]
\newtheorem{lemma}[theorem]{Lemma}
\newtheorem{proposition}[theorem]{Proposition}
\newtheorem{corollary}[theorem]{Corollary}
\newtheorem{definition}[theorem]{Definition}

%\newenvironment{proof}[1][Proof]{\begin{trivlist}
%\item[\hskip \labelsep {\bfseries #1}]}{\end{trivlist}}

\newenvironment{example}[1][Example]
{\begin{trivlist} \item[\hskip \labelsep {\bfseries #1}]}{\end{trivlist}}

\newenvironment{remark}[1][Remark]
{\begin{trivlist}\item[\hskip \labelsep {\bfseries #1}]}{\end{trivlist}}


\begin{document}
\sloppy 
\hbadness=10000
\vbadness=10000

\def\cal{\mathcal} % for LMCS
\def\calH{{\cal H}}

% journal, pa2

\def\PAtwoID{{\rm PA}_2{\rm ID}}
\def\CPAtwoID{{\rm CPA}_2{\rm ID}}
\def\AxiomPA{{\rm AxiomPA}}

% cmcs

\def\Co{{\rm co}}
\def\Axiom{{\rm Axiom}}
\def\Wk{{\rm Wk}}
\def\Cut{{\rm Cut}}
\def\BS{{\rm BS}}
\def\Bit{{\rm Bit}}
\def\LK{{\rm LK}}
\def\Head{{\rm head}}
\def\Tail{{\rm tail}}

% equivalence cyclic proofs

%\def\Uor{\displaystyle \mathop {\widetilde\bigvee}}
\def\Uor{\biguplus}
\def\CLKIDomega{{{\bf CLKID}^\omega}}
\def\CLKIDPA{{{\bf CLKID}^\omega+{\bf PA}}}
\def\LKIDPA{{{\bf LKID}+{\bf PA}}}
\def\PA{{\bf PA}}

% equiv note

\def\Choose{{\rm Choose}}
\def\Asym{{\rm Asym}}
\def\Peano{{\bf Peano}}
\def\Code#1{{\lceil{#1}\rceil}}
\def\Ineq{{\rm Ineq}}
\def\MonoPath{{\rm MonoPath}}
\def\MonoSeq{{\rm MonoSeq}}
\def\BoundedPath{{\rm BoundedPath}}
\def\InfPath{{\rm InfPath}}
\def\InfMonoPath{{\rm InfMonoPath}}
\def\InfMonoSeq{{\rm InfMonoSeq}}
\def\HomSeq{{\rm HomSeq}}
\def\InfHomSeq{{\rm InfHomSeq}}

\def\ErdosTree{{\rm ErdosTree}}
\def\KB{{\rm KB}}
\def\MonoList{{\rm MonoList}}
\def\Seq{{\rm Seq}}
\def\Univ{{\rm univ}}
\def\Infinite{{\rm Infinite}}
\def\S{{\bf S}}
\def\N{{\bf N}}
\def\Ind{{\rm Ind}}
\def\LKID{{\bf LKID}}
\def\PA{{\bf PA}}
\def\LKIDExt{{\bf LKIDExt}}
\def\CLKID{{\bf CLKID}}

% LICS

\def\Over#1#2{\deduce{#2}{#1}}
\def\List{{\rm List}}
\def\ListX{{\rm ListX}}
\def\ListO{{\rm ListO}}
\def\ListE{{\rm ListE}}
\def\Nl{{\rm nl}}

% CSLID

\def\SLLID{{\rm SLID}}
\def\SLID{{\rm SLID}}
\def\CSLID{{\rm CSLID}}
\def\CSLIDomega{{{\rm CSLID}\omega}}
\def\Eclass{{\rm Eclass}}
\def\Eq{{\rm Eq}}
\def\Deq{{\rm Deq}}
\def\Satom{{\rm Satom}}
\def\Cells{{\rm Cells}}
\def\Roots{{\rm Roots}}
\def\Elim{{\rm Elim}}
% \def\Case#1{{\rm Case-$#1$}}
\def\Case{{\rm Case}}
\def\Unfold{{\rm Unfold}}
\def\Pred{{\rm Pred}}
\def\Start{{\rm Start}}
\def\Unsat{{\rm Unsat}}
\def\Split{{\rm Split}}
\def\Underscore{\underline{\phantom{x}}}
\def\Rootcell{{\rm rootcell}}
\def\Rootshape{{\rm rootshape}}
\def\Jointleaf{{\rm jointleaf}}
\def\Jointleafimplicit{{\rm jointleafimplicit}}
\def\Jointnode{{\rm jointnode}}
\def\Directjoint{{\rm directjoint}}

% cyclicnote2

\def\Range{{\rm Range}}
\def\Xstart{X_{{\rm start}}}
\def\kstart{k_{{\rm start}}}
\def\Leaves{{\rm Leaves}}
\def\PureDist{{\rm PureDist}}
\def\Pure{{\rm Pure}}
\def\Identity{{\rm Identity}}
\def\Subst{{\rm Subst}}

% weak establish

\def\Connected{{\rm Connected}}
\def\Eststablished{{\rm Eststablished}}
\def\Valued{{\rm Valued}}

% monadic new

\def\SLRDbtw{{\rm SLRD}_{btw}}
\def\BTW{{{\rm BTW}}}
\def\Loc{{{\rm Loc}}}
\def\Ne{{\ne}}
\def\Nil{{{\rm nil}}}
\def\eqDef{=_{{\rm def}}}
\def\Inf#1{{\infty_{#1}}}
\def\Stores{{\rm Stores}}
\def\SVars{{{\rm SVars}}}
%\def\Val{{{\rm Val}}}
\def\MSO{{{\rm MSO}}}
\def\Sep{{{\rm Sep}}}
\def\Septwo{{{\rm SLMI}}}
\def\SLMI{{{\rm SLMI}}}
\def\Sepinf{{\rm Sep}\infty}
\def\THeaps{{\rm THeaps}}
\def\Root{{\rm Root}}
\def\TG{{\rm TGraph}}

% monadic

\def\nil{{{\rm nil}}}
\def\eqDef{=_{{\rm def}}}
\def\Inf#1{{\infty_{#1}}}
\def\Stores{{\rm Stores}}
\def\SVars{{{\rm SVars}}}
\def\Val{{{\mathcal V}}}
\def\MSO{{{\rm MSO}}}
\def\Sep{{{\rm sep}}}
\def\THeaps{{\rm THeaps}}
\def\Root{{\rm Root}}
\def\TG{{\rm TGraph}}

\def\Tilde{\widetilde}
\def\Bar{\overline}
\def\Lequiv{\Longleftrightarrow}
\def\Lto{\Longrightarrow}
\def\Lfrom{\Longleftarrow}
\def\Norm{{\rm Norm}}
\def\Noshare{{\rm Noshare}}
\def\Roots{{\rm Roots}}
\def\Forest{{\rm Forest}}
\def\Var{{\rm Var}}
\def\Range{{\rm Range}}
\def\Cell{{\rm Cell}}
\def\Tree{{\rm Tree}}
\def\Switch{{\rm Switch}}
\def\All{{\rm All}}
\def\To{\leadsto}
\def\tree{{\rm tree}}
\def\Paths{{\rm Paths}}
\def\Finite{{\rm Finite}}
\def\Const{{\rm Const}}
\def\Leaf{{\rm Leaf}}
\def\LeastElem{{\rm LeastElem}}
\def\LeastIndex{{\rm LeastIndex}}
\def\WSnS{{{\rm WSnS}}}
\def\Expand{{{\rm Expand}}}

% previous

\long\def\J#1{} % skip
\def\T#1{\hbox{\color{green}{$\clubsuit #1$}}}
\def\W#1{\hbox{\color{Orange}{$\spadesuit #1$}}}

\def\Node{{{\rm Node}}}
\def\LL{{{\rm LL}}}
\def\DSN{{{\rm DSN}}}
\def\DCL{{{\rm DCL}}}
\def\LS{{{\rm LS}}}
\def\Ls{{{\rm ls}}}

\def\FPV{{{\rm FPV}}}
\def\Lfp{{{\rm lfp}}}
\def\IsHeap{{{\rm IsHeap}}}

\def\Equiv{\quad \equiv\quad }
\def\Null{{{\rm null}}}
\def\Emp{{{\rm emp}}}
\def\If{{{\rm if\ }}}
\def\Then{{{\rm \ then\ }}}
\def\Else{{{\rm \ else\ }}}
\def\While{{{\rm while\ }}}
\def\Do{{{\rm \ do\ }}}
\def\Cons{{{\rm cons}}}
\def\Dispose{{{\rm dispose}}}
\def\Vars{{{\rm Vars}}}
\def\Locs{{{\rm Locs}}}
\def\States{{{\rm States}}}
\def\Heaps{{{\rm Heaps}}}
\def\FV{{{\rm FV}}}
\def\True{{{\rm true}}}
\def\False{{{\rm false}}}
\def\Dom{{{\rm Dom}}}
\def\Abort{{{\rm abort}}}
\def\New{{{\rm New}}}
\def\W{{{\rm W}}}
\def\Pair{{{\rm Pair}}}
\def\Lh{{{\rm Lh}}}
\def\lh{{{\rm lh}}}
\def\Elem{{{\rm Elem}}}
\def\EEval{{{\rm EEval}}}
% \def\BEval{{{\rm BEval}}} % not used 
\def\PEval{{{\rm PEval}}}
\def\HEval{{{\rm HEval}}}
\def\EVal{{{\rm Eval}}}
\def\Domain{{{\rm Domain}}}
\def\Exec{{{\rm Exec}}}
\def\Store{{{\rm Store}}}
\def\Heap{{{\rm Heap}}}
\def\Storecode{{{\rm Storecode}}}
\def\Heapcode{{{\rm Heapcode}}}
\def\Lesslh{{{\rm Lesslh}}}
\def\Addseq{{{\rm Addseq}}}
\def\Separate{{{\rm Separate}}}
\def\Result{{{\rm Result}}}
\def\Lookup{{{\rm Lookup}}}
\def\ChangeStore{{{\rm ChangeStore}}}
\def\ChangeHeap{{{\rm ChangeHeap}}}
\def\Wand{\mathbin{\hbox{\hbox{---}$*$}}}
\def\Eval#1{[\kern -1pt[{#1}]\kern -1pt]}
\def\Vec{\overrightarrow}
%\def\Vec#1{{\bf #1}}

\def\Tilde{\widetilde}
\def\Break{\hfil\break\hbox{}}

\newcommand{\freshcopy}           [1]                      {  {\tilde{#1}} }
%\newcommand{\freshcopy}           [1]                      {  \displaystyle{\mathop{#1}^{\sim}} }
\newcommand{\bfColor} [2]     {{\bf \color{#1}{#2}}}
\newcommand{\arity}               {{\tt ar}}          
\newcommand{\trace}               {{\tt tr}} 
\newcommand{\stat}                 {{\tt stat}} 
\newcommand{\Stat}                {{\tt S}}     
\newcommand{\prog}               {{\tt p}} 
\newcommand{\Prog}               {{\tt P}}   
\newcommand{\fun}                 {{\tt fun}}          
\newcommand{\quotationMarks} [1]  {{\textquotedblleft #1\textquotedblright}}
\newcommand{\universe} [1]  {{| #1 |}}
\newcommand{\Rule}               {{\tt Rule}} 
\newcommand{\Path}               {{\tt Path}} 
\newcommand{\pathBetween} {{\tt path}} 
\newcommand{\Trace}               {{\tt Trace}} 
\newcommand{\Cyclic}              {{\tt Cyclic}}
\newcommand{\Cycle}              {{\tt Cycle}}
\newcommand{\dom}              {{\tt dom}}
\newcommand{\codom}              {{\tt codom}}
\newcommand{\Periodic}              {{\tt Periodic}}
\newcommand{\setR}              {{\mathcal R}}
\newcommand{\setL}              {{\mathcal L}}
\newcommand{\setG}              {{\mathcal G}}
\newcommand{\Pfin}      {{ {\mathcal R}_{\mbox{\tiny fin}} }}
\newcommand{\Card}              {{\tt Card}}
\newcommand{\restr}               {{\lceil}}
\newcommand{\comp}              {{\circ}}
\newcommand{\val}                   {{\tt val}}
\newcommand{\setP}                   {{\mathcal P}}
\newcommand{\Rprog}               {{\star}}
\newcommand{\RProg}               {{\bigstar}}
\newcommand{\length}               {{\mbox{\tt len}}}
\newcommand{\lex}                    {{\mbox{\tiny lex}}}
\newcommand{\setH}                  {{\mathcal H}}
\newcommand{\setS}                  {{\mathcal S}}
\newcommand{\setT}                  {{\mathcal T}}
\newcommand{\id}                      {{\tt id}}
\newcommand{\BrowerKleene}  {{\mbox{\tiny BK}}}
\newcommand{\conc}                 {{*}}




% MORE MACRO
\newcommand{\size} 		                  {{\tt s}}
\newcommand{\state} 		               {{\tt S}}
\newcommand{\minStable}             {{\tt m}}
\newcommand{\unfold}                   {{\tt u}}
\newcommand{\height}                    {{\tt h}}
\newcommand{\setF} 		                   {{\mathcal F}}
\newcommand{\setU} 	                   {{\mathcal U}}
\newcommand{\cons}                       {{\tt cons}}
\newcommand{\newton} [2]             {{\genfrac(){0pt}{1} {#1}{#2}}}
\newcommand{\questionMarks} [1] {{‘‘#1’’}}
\newcommand{\DomCard}               {{\tt DomCard}}
\newcommand{\MaxInd}                  {{\tt M}}
%END MACRO

\title{A Poly-exponential Translation from Injective Cyclic to Inductive Proofs}
\date{}

\author[S. Berardi]{Stefano Berardi}
\address{Universit\`{a} di Torino,
Torino, Italy}
\email{stefano@unito.it}

%% mandatory lists of keywords 
\keywords{
proof theory
inductive definitions
cyclic proofs
infinite Ramsey theorem
}

%%%%%%%%%%%%%%%%%%%%%%%%%%%%%%%
% Titolo originario: An Intuitionistic and First Order proof that GTC implies Correctnes
%%%%%%%%%%%%%%%%%%%%%%%%%%%%%%%
% questo lavoro mio in data 16/07/2025 è diventato un lavoro su injectivity e prove
% cicliche, per CSL 2026.  Ci sono ora tre nomi, io, Osiero e de' Liguoro. Viene rifiutato
% in data 15 ottobre 2026. Un recensore, forse Amram, è stato particolarmente
% critico, sostiene che il lavoro non contiene alcun progresso rispetto al suo
% lavoro SCT con Niel Jones. In data 22 Ottobre provo a riprendere il lavoro del 
% 16/072025 per vedere se si può trasformare in un lavoro di proof-theory e 
% trasformazione di prove cicliche in prove ordinarie, che non possa venir recensito
% da amram.
%%%%%%%%%%%%%%%%%%%%%%%%%%%%%%%


\maketitle


\emph{Proof started in Torino, January 20, 2026. 
Ugo suggested adding complexity results to motivate the new algorithm}.

Mirai \cite{} provided a translation from from cyclic to inductive proofs. The proof builds
on an existing algorithm for generating ranking functions and 
requiring triple exponential time. A first improvement would be to switch to an equivalent
formulation of cyclic proofs with injective bicolored relation (i.e. fan-in free, with at most one
arrow entering in each node). In this case, the same algorithm provides a double exponential
time \cite{} and an inductive proof in output of size $O(2^{n \log(n)})$.
In this paper we further improve the translation to a single exponential (poly-exponential) time, 
by introducing a new algorithm. The price is having an inductive proof in output
of poly-exponential size, instead than an ouput of size $O(2^{n \log(n)})$.


\newpage   \section{Quick Overview of the GTC and PDC Conditions}
Let us start from a cyclic proof $\Pi$ with set $\setR$ of bicolored relations on a set 
$[1,\MaxInd]$. $\MaxInd$ is the maximum number of inductive formula in the left-hand side 
of some sequent in $\Pi$. Each relation $R$ is associated to an arrow $a \rightarrow b$ 
between nodes of $\Pi$, and any two relations $R:a \rightarrow b$, $S:b \rightarrow c$
are said to be composable.
 
Assume $\setR^+$ is the closure of $\setR$ under composition. The Progressing Domain 
Condition (PDC) for a bicolored relation $R : a \rightarrow a$ says: 
there are $x, y \in \dom(R)$ such that $R$ is progressing on $(x,y)$.
We say that set $\setR$ of bicolored relations on $\Pi$
satisfies PDC if and only if all $R \in \setR$ with $R: a \rightarrow a$ satisfy PDC. 
Remark that the Global Trace Condition (GTC) is a property of $\setR$ as a whole, 
while PDC is a property of single relations $R \in \setR$, $R:a \rightarrow a$. 

We assume that:

\begin{enumerate}
\item  
$\setR$ is a set of injective relations, hence $\setR^+$ is a set of injective relations, too.
\item
$\setR^+$ satisfies PDC 
\end{enumerate}

PDC for injective relations is equivalent to GTC, in general GTC is a stronger condition
than PDC. 
\\

{\bf An example.} Let $R = (\{(1,1),(2,2)\},\{(2,1)\})$ with $R$ progressing
on $(2,1)$. Then $R$ satisfies PDC on the graph $\{(a,a)\}$, and $RR=R$, 
therefore  $\setR^+=\{R\}$ and $\setR^+$ satisfies PDC. However, the only traces any
infinite path are $(1,1,1,\ldots)$ and $(2,2,2,\ldots)$ and $(2,2,2,\ldots,1,1,1,\ldots)$. The
first two traces never progress, the third trace progresses once, on $(2,1)$. So $\setR^+$
does \emph{not} satisfy GTC. The fact that GTC is false for $\setR^+$ 
is not in contradiction with what we said, because $R \in \setR^+$ is  \emph{not} injective.
\\

Suppose we have a path $\pi$ from a companion c to itself, moving back  through $c$
for $k \ge 1$ times: we have $c \rightarrow c \rightarrow \ldots \rightarrow c$, with $k$ arrows.
For any $i \in \{1, \ldots, k\}$, suppose the step $c \rightarrow c$ having number $i$
is associated to some bicolored relation $R_i : c \rightarrow c$, $R_i \in \setR^+$. 
The path $\pi$ itself is associated to the sequence $\vec{R} = R_1 \ldots R_k$ of such relations. 
For the translation, we have to associate a measure to each such sequence $\vec{R}$
which is decreasing by extension.

In order to do so, we first to derive some elementar properties of the sequences of uncolored injective relations on $[1,\MaxInd]$. This properties having nothing to do with colors and
with cyclic proofs.


\newpage   \section{Domain-Stable Sequences and Grouping}
From now on we write $\N$ for the set of non-negative integers and for any $a,b \in \N$
we write $[a,b] = \{c \in \N | a \le c \le b\}$. 
We consider binary relations on a finite set: for sake of simplicity we only 
consider the finite sets $[1,\MaxInd]$, for some $\MaxInd \in \N$ which we assume to be
fixed. We will abbreviate \questionMarks{binary relations on $[1,\MaxInd]$} with
\questionMarks{relations}. 

\begin{definition}[Grouping]
A grouping of a sequence $\vec{S} = S_1, \ldots, S_n$ of relations is a sequence
of non-empty products of relations, defined by induction on $n$ as follows.
\begin{enumerate}
\item
If $n=0$, then $\vec{S}$ (empty) is the unique grouping of $\vec{S}$.
\item
Suppose $T_1, \ldots, T_m$ is any grouping of some $\vec{S}$. 
Then $T_1, \ldots, T_{i-1}, (T_i \ldots T_m S)$ is a grouping of 
$\vec{S},S$, for all $i \in [1,m]$. 
\end{enumerate}
\end{definition}

To put otherwise, the groupings of $S_1, \ldots, S_n$ are all and only 
sequences of products of the form: 
$$
S_{a_1} \ldots S_{b_1}, \ \ \ 
S_{a_2} \ldots S_{b_2}, \ \ \ 
\ldots, \ \ \ 
S_{a_{h+1}} \ldots S_{b_{h+1}}
$$
The segments $[a_1,b_1], \ldots, [a_{h+1},b_{h+1}]$ are a decomposition
of $[1,n]$ into a disjoint union for some increasing sequence 
$0 < b_1 < b_2 < i_3 < \ldots < b_h < b_{h+1} = n$. We necessarily have
$a_0=1$ and $a_{i+1} = b_i+1$ for all $i \in [1,h]$. 

Some trivial examples. 
\begin{enumerate}
\item
We can have $h=0$, $b_1=n$ and the trivial grouping consisting of a single product  
$S_1 \ldots S_n$. 
\item
We can have $h=m$, $b_1=1$, \ldots, $b_m=n$ and the trivial grouping consisting of the 
sequence $S_1, \ldots, S_n$ itself, considered a sequence of unary products.
\end{enumerate}

We first define some domain inclusions and equalities 
which are valid for all sequences of relations, then we isolate the \questionMarks{stable}
grouping, in which more domain inclusions are valid. 

\begin{lemma}[Domain-inclusion]
\label{lemma-domain-inclusion}
Assume $R, S, T$ are any relations and $R_1, \ldots, R_n$ is any sequence of relations.
\begin{enumerate}

\item
\label{lemma-domain-inclusion-01}
For any two relations $R$, $S$ we have $\dom(RS) \subseteq \dom(R)$
and therefore $\Card(\dom(RS)) \le \Card(\dom(R))$.

\item
\label{lemma-domain-inclusion-02}

$$
\dom(R_1)   \ \ \  \supseteq  \ \ \ 
\dom(R_1 R_2)    \ \ \  \supseteq  \ \ \ 
\dom(R_1 R_2 R_3)    \ \ \  \supseteq  \ \ \  
\ldots   \ \ \  \supseteq
\dom(R_1 \ldots R_n)
$$

\item
\label{lemma-domain-inclusion-03}
$
(\dom(R) = \dom(RS))  
\ \ \  \Leftrightarrow \ \ \ 
(\dom(R) \subseteq \dom(RS)) 
\ \ \  \Leftrightarrow \ \ \  
(\dom(R) \not \supset  \dom(RS))
\ \ \  \Leftrightarrow \ \ \ 
(\Card(\dom(RS)) = \Card(\dom(R))
$.

\item
\label{lemma-domain-inclusion-04}
If $\dom(S) = \dom(ST)$ then 
$\dom(RS) = \dom(RST)$.

\end{enumerate}
\end{lemma}

\begin{proof}
\begin{enumerate}

\item
%\label{lemma-domain-inclusion-01}
We have to prove that $\dom(RS) \subseteq \dom(R)$. For, we assume $x \in \dom(RS)$ in 
order to prove that $x \in \dom(R)$. If $x \in \dom(RS)$ then for some $z$ we have $x RS z$, 
therefore $x R y S z$ for some $y$. We conclude that $x \in \dom(R)$, as wished.

\item
%\label{lemma-domain-inclusion-02}
We prove 
$
\dom(R_1) \supseteq 
\dom(R_1 R_2) \supseteq 
\dom(R_1 R_2 R_3) \supseteq 
\ldots \supseteq 
\dom(R_1 \ldots R_n)
$ 
by induction on $n$. 
When $n=0,1$ there is no inclusion to prove. Suppose $n \ge 1$ and the thesis 
for $n$, namely that $\dom(R_1) \supseteq \dom(R_1 R_2) \supseteq \dom(R_1 R_2 R_3) 
\supseteq \ldots \supseteq \dom(R_1 \ldots R_n)$. We have to prove the thesis for $n+1$. 
By point \ref{lemma-domain-inclusion-01}above 
we have $\dom(R_1 \ldots R_{n} R_{n+1}) \subseteq \dom(R_1 \ldots R_{n})$, namely. 
$\dom(R_1 \ldots R_{n}) \supseteq \dom(R_1 \ldots R_{n} R_{n+1})$. 
From this and the inductive hypothesis we deduce the thesis for $n+1$.

\item
%\label{lemma-domain-inclusion-03}
If $(\dom(R) = \dom(RS))$ then $(\dom(R) \subseteq \dom(RS))$. If
$(\dom(R) \subseteq \dom(RS))$, then $(\dom(R) \not \supset \dom(RS))$.
If $(\dom(R) \not \supset \dom(RS))$, then from $(\dom(R) \supseteq \dom(RS))$
(point \ref{lemma-domain-inclusion-01}above) we deduce that $(\dom(R) = \dom(RS))$,
hence $(\Card(\dom(R)) = \Card(\dom(RS)))$. If $(\Card(\dom(R)) = \Card(\dom(RS)))$
then from $(\dom(R) \supseteq \dom(RS))$ and the fact that both sets are included in
$[1,\MaxInd]$ and therefore finite we deduce that $(\dom(R) = \dom(RS))$.

\item
%\label{lemma-domain-inclusion-04}
Assume that $\dom(S) = \dom(ST)$ in order to prove that 
$\dom(RS) = \dom(RST)$. 
We have $\dom(RS) \supseteq \dom(RST)$ by point \ref{lemma-domain-inclusion-01} above,
so we only have to prove that $\dom(RS) \subseteq \dom(RST)$.
For, we assume $a \in \dom(RS)$ in order to prove that $a \in \dom(RST)$. 
From $a \in \dom(RS)$ we deduce that there are $b \in \dom(S)$ and $c$ such that $a R b S c$. 
From  $\dom(S) = \dom(ST)$ we deduce that $b \in \dom(ST)$, namely that
for some $c' \in \dom(T)$ and $d$ we have $b S c' T d$. From $a R b S c' T d$ we conclude
that $a \in \dom(RST)$, as wished.

\end{enumerate}
\end{proof}

The inclusion $\dom(RS) \subseteq \dom(R)$ proved in
Lemma \ref{lemma-domain-inclusion}.\ref{lemma-domain-inclusion-01}
is often strict. An example: if $R = \{(1,1)\}$ and $S = \{(2,2)\}$ then $RS=\emptyset$,
therefore $\dom(RS) = \emptyset \subset \dom(R) = \{1\}$. 
We call domain-stable any sequence $\vec{R}$ of relations in which this never happens.
 
\begin{definition}[Domain-Stability]
Let $R$, $S$ be any two relations and $\vec{S}=S_1, \ldots, S_n$ be any sequence of relations. 
\begin{enumerate}

\item
We say that a pair $R, S$ is \emph{domain-stable}, only \emph{stable} for short, if 
and only if $\dom(R) = \dom(RS)$.

\item
$\vec{S}$ is stable if and only if for all $i \in [1,n-1]$
the pair $S_i$, $S_{i+1}$ is stable.

\end{enumerate} 
\end{definition}


A sequence of length $n=0,1$ is trivially stable because $[1,n-1] = \emptyset$.
By Lemma \ref{lemma-domain-inclusion}.\ref{lemma-domain-inclusion-03} we can
restate domain-stability in several ways. We have
$R,S$ stable $ \Leftrightarrow (\dom(R) = \dom(RS))  
  \Leftrightarrow  
(\dom(R) \subseteq \dom(RS)) 
  \Leftrightarrow   
(\dom(R) \not \supset  \dom(RS))
$.

Another way of saying that the pair $R,S$ is stable is that we do not remove elements
from $\dom(R)$ when we compose $R$ with $S$.

In particular, $R,S$ stable means that $(\dom(R) \subseteq \dom(RS))$, 
namely that whenever $x$ has an $R$-image $y$, then $x$ has some
$R$-image $y'$ having some $S$-image $z$. It does \emph{not} say, however,
that each $R$-image of $x$ has some $S$-image: it only says that if $x$ has some
$R$-image, then some of the $R$-images of $x$ have some $S$-image. 

In the same way we can restate the stability for a sequence.

\begin{proposition}[Restating Stability]
Let $\vec{S}=S_1, \ldots, S_n$ be any sequence of relations. 
Then  $\vec{S}$ is stable if and only if $S_i \restr \dom(S_i) \times \dom(S_{i+1})$ 
is everywhere defined for all $i \in [1,n-1]$. 
\end{proposition}

\begin{proof}
Indeed, this is just another way of saying that if $x$ has some $S_i$-image,
then some $S_i$-image of $x$ has some $S_{i+1}$-image. This is the definition
of $S_i$, $S_{i+1}$ domain-stable for $i \in [1,n-1]$. 
\end{proof}

Domain-stable sequences satisfy more domain equalities than generic sequences.


\begin{lemma}[Domain-Stability]
\label{lemma-domain-stability}
Let $\vec{S}=S_1, \ldots, S_n$ be any sequence of relations. The following are equivalent:

\begin{enumerate}
\item
\label{lemma-domain-stability-01}
$\vec{S}$ is stable

\item
\label{lemma-domain-stability-02}
for all $i, j \in [1, n]$, $i \le j$ we have:
$$
\dom(S_i) =  \dom(S_i \ldots S_j)
$$
or equivalently $\Card(\dom(S_i)) =  \Card(\dom(S_i \ldots S_j))$.

\item
\label{lemma-domain-stability-03}
all grouping and all prefixes of $\vec{S}$ are stable.

\item 
\label{lemma-domain-stability-04}
for all $i \in [1,n-1]$ we have:
$$
\dom(S_i) = \dom(S_i \ldots S_n)
$$

\end{enumerate}



\end{lemma}

\begin{proof}
We can equivalently restate $\dom(S_i) =  \dom(S_i \ldots S_j)$ as
$\Card(\dom(S_i)) =  \Card(\dom(S_i \ldots S_j))$, 
given that we have $\dom(S_i) \supseteq  \dom(S_i \ldots S_j)$ 
by Lemma \ref{lemma-domain-inclusion}.\ref{lemma-domain-inclusion-02})
and the fact that they are subsets of the finite set $[1,\MaxInd]$.
Therefore we can skip the formulation using 
$\Card(\dom(S_i)) =  \Card(\dom(S_i \ldots S_j))$
when deriving the equivalences about point \ref{lemma-domain-stability-02}.

We will now prove the equivalence of points 
\ref{lemma-domain-stability-01},
\ref{lemma-domain-stability-02},
\ref{lemma-domain-stability-03},
namely we prove that
\begin{center}
\ref{lemma-domain-stability-01} 
$\implies$ 
\ref{lemma-domain-stability-02}
$\implies$ 
\ref{lemma-domain-stability-03} 
$\implies$ 
\ref{lemma-domain-stability-01}
\end{center}
Then we prove the equivalence of points
\ref{lemma-domain-stability-02},
\ref{lemma-domain-stability-04},
by proving that 
\begin{center}
\ref{lemma-domain-stability-02} $\implies$  
\ref{lemma-domain-stability-04} $\implies$  
\ref{lemma-domain-stability-02}
\end{center}


\begin{enumerate}

\item
\ref{lemma-domain-stability-01} $\implies$ \ref{lemma-domain-stability-02}.

We assume point \ref{lemma-domain-stability-01}: $\vec{S}$ is stable, in to prove 
point \ref{lemma-domain-stability-02}: 
for all $i, j \in [1, n]$, $i \le j$ we have $\dom(S_i) =  \dom(S_i \ldots S_j)$.

For, we assume $j \ge i$ and we prove $\dom(S_i) =  \dom(S_i \ldots S_j)$ by induction on $j$.
\begin{enumerate}
\item
Assume $j=i$. Then $\dom(S_i) =  \dom(S_i)$.
\item
Assume $\dom(S_i) =  \dom(S_i \ldots S_j)$ and $j<n$ in order to prove that 
$\dom(S_i) =  \dom(S_i \ldots S_j S_{j+1})$. By defiition of stability we have 
$\dom(S_j) = \dom(S_j S_{j+1})$. By Lemma 
\ref{lemma-domain-inclusion}.\ref{lemma-domain-inclusion-04} 
applied to $(S_i \ldots S_{j-1}), \ S_j, \ S_{j+1}$ we deduce 
$\dom( (S_i \ldots S_{j-1}) S_j) = \dom( (S_i \ldots S_{j-1}) S_j S_{j+1})$.
From the inductive hypothesis $\dom(S_i) =  \dom( (S_i \ldots S_{j-1}) S_j)$
we conclude that $\dom(S_i) =  \dom((S_i \ldots S_{j-1}) S_j S_{j+1})$, as wished.
\end{enumerate}


\item
\ref{lemma-domain-stability-02} $\implies$ \ref{lemma-domain-stability-03}.

We assume point \ref{lemma-domain-stability-02}: 
for all $i, j \in [1, n]$, $i \le j$ we have $\dom(S_i)  =  \dom(S_i \ldots S_j)$,
in order to prove point \ref{lemma-domain-stability-03}:
all prefixes and grouping of $\vec{S}$ are stable. 

All prefixes are stable because the finite quantification defining stability for a prefix 
is a restriction of the finite quantification defining stability for the sequence. 

We prove that all grouping of $\vec{S}$ are stable. We assume that
$i, j,k\in [1, n]$, $i \le j < k$ and that $S_i \ldots S_j$, $S_{j+1} \ldots S_k$
are two consecutive products in some grouping of $\vec{S}$, and we have to prove 
that $\dom(S_i \ldots S_j) = \dom( (S_i \ldots S_j) (S_{j+1} \ldots S_k))$. 
By point \ref{lemma-domain-stability-02} we have 
$$
\dom(S_i \ldots S_j) = \dom(S_i) = \dom(S_i \ldots S_j S_{j+1} \ldots S_k)
$$
This proves point \ref{lemma-domain-stability-03}.


\item
\ref{lemma-domain-stability-03} $\implies$ \ref{lemma-domain-stability-01}.

We assume point \ref{lemma-domain-stability-03}: all prefixes 
and grouping of $\vec{S}$ are stable,
and we have to prove point \ref{lemma-domain-stability-01}: $\vec{S}$ is stable. 

This follows from the fact that $\vec{S}$ itself is a trivial case of grouping of $\vec{S}$.


\item
\ref{lemma-domain-stability-02} $\implies$  \ref{lemma-domain-stability-04}.
We assume point \ref{lemma-domain-stability-02}:
for all $i, j \in [1, n]$, $i \le j$ we have $\dom(S_i) =  \dom(S_i \ldots S_j)$
and we have to prove point \ref{lemma-domain-stability-04}: 
for all $i \in [1, n-1]$ we have $\dom(S_i) =  \dom(S_i \ldots S_n)$.

This follows by choosing $j=n \ge i$.


\item
\ref{lemma-domain-stability-04} $\implies$  \ref{lemma-domain-stability-02}.

We assume point \ref{lemma-domain-stability-04}: 
for all $i \in [1, n-1]$ we have $\dom(S_i) =  \dom(S_i \ldots S_n)$, and we have to prove
point \ref{lemma-domain-stability-02}: 
for all $i, j \in [1, n]$, $i \le j$ we have $\dom(S_i) =  \dom(S_i \ldots S_j)$.

\emph{Suppose $i=n$}. 
Then the thesis becomes $\dom(S_n) =  \dom(S_n)$, which is trivially true.

\emph{Suppose $i \in [1,n-1]$}.
By Lemma \ref{lemma-domain-inclusion}.\ref{lemma-domain-inclusion-01},
for each sequence $\vec{S}$ of relations we have 
$$
\dom(S_i)  \ \ \ \supseteq \ \ \ 
\dom(S_iS_{i+1})  \ \ \  \supseteq \ \ \ 
\dom(S_iS_{i+1}S_{i+2})  \ \ \  \supseteq  \ \ \ 
\ldots  \ \ \  \supseteq  \ \ \ 
\dom(S_i \ldots S_n)
$$ 
therefore if we add the assumption $\dom(S_i \ldots S_n) = \dom(S_i)$ we conclude that
$$
\dom(S_i)   \ \ \ =  \ \ \ 
\dom(S_iS_{i+1})  \ \ \  =  \ \ \ 
\dom(S_iS_{i+1}S_{i+2})   \ \ \ =  \ \ \ 
\ldots   \ \ \ =  \ \ \ 
\dom(S_i \ldots S_n)
$$ 
as wished.
 
\end{enumerate}
\end{proof}

Lemma \ref{lemma-domain-stability}.\ref{lemma-domain-stability-02}
says that a sequence is stable if and only if whenever $x$ has some $S_i$-image,
then $x$ has some chain of $S_i, S_{i+1}, S_{i+2}, \ldots$-images of any length
moving through the sequence. More informally, that we never drop any element of $\dom(S_i)$
when composing with $S_{i+1}$, $S_{i+2}$, \ldots.
This is an alternative definition of domain-stable sequence.

We compare two groupings, looking for the grouping whose product closet to the left 
are shorter.

\begin{definition}[Lexicographic Ordering on Grouping]
\label{definition-grouping-order}
Assume $\vec{S}=S_1, \ldots, S_n$ is a sequence of relations and 
$$
\vec{R} \ \ \ = \ \ \ 
S_{a_1} \ldots S_{b_1}, \ \ \ S_{a_2} \ldots S_{b_2}, \ \ \ \ldots, \ \ \ S_{a_{h+1}} \ldots S_{b_{h+1}}
$$
is a grouping of $\vec{S}$ with $a_1=1, a_{i+1} = b_i+1$ for all $i \in [1,h]$ and 
$b_{h+1} = n$.

\begin{enumerate}
\item
The degree of the grouping $\vec{R}$ of $\vec{S}$ is the sequence
of the lengths of its products, namely $(b_1, b_2 - b_1, \ldots, b_{h+1} - b_h)$. 
\item
We compare two grouping of the same sequence by comparing their degrees in the 
lexicographic ordering.
\end{enumerate}
\end{definition} 

For instance, 
$S_1 \ldots S_n$ is the maximum grouping and $S_1, \ldots, S_n$ the minimum grouping.

We use the ordering on grouping to select a simplest stable grouping of $\vec{S}$.
All sequences $\vec{S}=S_1, \ldots, S_n$ have a trivial stable grouping, by taking the
product $S_1 \ldots S_n$ of the sequence. This is the maximum grouping, 
and therefore it is the maximum stable grouping, but it is trivial.

Lexicographic ordering is total, therefore there is a minimum stable grouping.

\begin{definition}[The Minimum Stable Grouping]
\label{definition-minstable}
Assume $\vec{S}$ is any sequence of relations. We define 
$\minStable(\vec{S})$ as the minimum stable grouping of $\vec{S}$ in the
lexicographic ordering.
\end{definition}


We provide an alternative characterization of the minimum stable grouping of $\vec{S}$
using a map $\phi(\vec{S},\cdot) : [1,\MaxInd] \rightarrow  [1,\MaxInd]$. 
Then using the map  map $\phi(\vec{S},\cdot)$ we will provide an alternative definition of 
$\minStable(\vec{S})$ by induction on $\vec{S}$, which is computable in polynomial time. 
We first need the notion of matching between $i,j \in [1,n]$.

\begin{definition}[Stable Matching]
Let $\vec{S} = S_1, \ldots, S_n$ and $i,j \in [1,n]$, $i \le j$.
We say that $j$ is a stable matching of $i$ in $\vec{S}$
if and only if $\dom(S_i \ldots S_j) = \dom(S_i \ldots S_n)$.
\end{definition}

By definition unfolding, $j$ is a stable matching of $i$ in $\vec{S}$
if and only if all $x \in [1,\MaxInd]$ having some $S_i \ldots S_j$-image have some
$S_i \ldots S_n$-image and conversely. 

Thus, $j=n$ is trivially a stable matching of any $i \in [1,n]$.

\begin{definition}[Minimum Stable Matching $\phi(\vec{S},i)$]
Let $\vec{S} = S_1, \ldots, S_n$ and $i \in [1,n]$.
We define the integer $\phi(\vec{S},i) \in [i,n]$ as the minimum stable matching $j\in [i,n]$ 
of $i$ in $\vec{S}$. We write just $\phi(i)$ when $\vec{S}$ is clear from the context,
\end{definition}

We define a stable grouping $\vec{S}^{\phi}$ of $\vec{S}$ using the map $\phi$.

\begin{definition} [A Stable Grouping Defined by the Map $\phi$]
$\vec{S}^{\phi}=R_1, \ldots, R_{h+1}$ and: 
$$
R_1 = (S_{a_1}\ldots S_{\phi(a_1)}), \ \ \ 
R_2 = (S_{a_2}\ldots S_{\phi(a_2)}), \ \ \ 
\ldots, \ \ \ 
R_{h+1} = (S_{a_{h+1}}\ldots S_{\phi(a_{h+1})}
$$ 
and $a_1=1$ and $a_{i+1} = \phi(a_i)+1$ 
for all $i \in [1,h]$ and $\phi(a_{h+1}) = n$. 
\end{definition}

\begin{lemma} [A Stable Grouping Defined by the Map $\phi$]
\label{lemma-stable-phi-1}
$\vec{S}^\phi$ is a stable grouping of $\vec{S}$.
\end{lemma}

\begin{proof}
Let $\vec{R} = \vec{S}^\phi$. We have
$
\dom(R_i) = 
\dom(S_{a_i}\ldots S_{\phi(a_i)}) = 
$ 
(by definition of $\phi$)  = 
$
\dom(S_{a_i}\ldots S_{n}) =
\dom(R_{i}\ldots R_h)
$
We deduce that the grouping $\vec{R}$ is stable by 
Lemma \ref{lemma-domain-stability}.\ref{lemma-domain-stability-04}.
\end{proof}

$\vec{S}^\phi$ is the minimum stable grouping.

\begin{lemma} [The Minimum Stable Grouping is Definable by the Map $\phi$]
\label{lemma-stable-phi-2}
Stable groupings of $\vec{S}$ are exactly the grouping of $\vec{S}^\phi$
and $\vec{S}^\phi=\minStable{\vec{S}}$, the minimum stable grouping of $\vec{S}$.
\end{lemma}


\begin{proof}
We proved in \ref{lemma-stable-phi-1} that $\vec{S}^\phi$ is stable. We have still to
prove that all stable grouping of $\vec{S}$ 
are less or equal than $\vec{S}^\phi$ in the lexicographic ordering.

We will in fact prove that stable grouping of $\vec{S}$ are exactly the grouping of 
$\vec{S}^\phi$, we will deduce that they are all less or equal than $\vec{S}^\phi$.

From $\vec{S}$ stable and 
Lemma \ref{lemma-domain-stability}.\ref{lemma-domain-stability-03} 
we deduce that all grouping of $\vec{S}$ are stable.

Conversely, we have to prove that any stable grouping of $\vec{S}$  is a grouping
of $\vec{S}^\phi$.

We first check that in any stable grouping of $\vec{S}$ any 
product $S_{a_i+1} \ldots S_{a_{i+1}}$  must include $S_{\phi(a_i+1)}$.

Indeed, stability of the grouping and 
Lemma \ref{lemma-domain-stability}.\ref{lemma-domain-stability-04} 
imply
$$
\dom(S_{a_i+1} \ldots S_{a_{i+1}}) = 
\dom(S_{a_i+1} \ldots S_{a_{i+1}} \ \ \ S_{a_{i+1}+1} \ldots S_{a_{i+2}} \ \ \ \ldots) =
\dom(S_{a_i+1} \ldots S_n)
$$ 
Then by definition of $\phi$ we have $\phi(a_i+1) \le a_{i+1} $.

Therefore in any stable grouping of $\vec{S}$
any product $S_i \ldots S_j$ is either $(S_{i}\ldots 
S_{\phi(i)})$, or if $\phi(i) < j$ it includes $(S_{\phi(i)+1}\ldots S_{\phi(\phi(i)+1)})$,
and so forth. By induction on $j-i$ we can prove that $S_i \ldots S_j$ has the form
$(S_{a_1}\ldots S_{\phi(a_1)}) (S_{a_2}\ldots S_{\phi(a_2)})
\ldots S_{a_{k+1}}\ldots S_{\phi(a_{k+1})}$ with $a_1=i$ and $a_{b+1} = \phi(a_b)+1$ 
for all $b \in [1,k]$ and $\phi(a_{k+1}) = j$. 


We conclude that all stable grouping are greater or equal than the following grouping 
$\vec{S}^\phi =R_1, \ldots, R_{h+1}$: 
$$
R_1 = (S_{a_1}\ldots S_{\phi(a_1)}), \ \ \ 
R_2 = (S_{a_2}\ldots S_{\phi(a_2)}), \ \ \ 
\ldots, \ \ \ 
R_{h+1} = (S_{a_{h+1}}\ldots S_{\phi(a_{h+1})}
$$ 
where $a_1=1$ and $a_{i+1} = \phi(a_i)+1$ 
for all $i \in [1,h]$ and $\phi(a_{h+1}) = n$. 
\end{proof}

$\phi(\vec{S},i)$ can be computed by induction on $\vec{S}$, as follows. 

\begin{lemma}[Characterization of $\phi(\vec{S},i)$]
\label{lemma-phi-characterization}
We have $\phi(\vec{S} S_{n+1}, n+1) = n+1$ and for all $i \in [1,n]$
one of the following holds:
\begin{enumerate}
\item
either $\dom(S_i \ldots S_n) = \dom(S_i \ldots S_n S_{n+1})$
and $\phi(\vec{S}S_{n+1},i) = \phi(\vec{S},i)$;
\item
or $\dom(S_i \ldots S_n) \supset \dom(S_i \ldots S_n S_{n+1})$
and $\phi(\vec{S}S_{n+1},i) = n+1$.
\end{enumerate}
$\phi(\vec{S},\cdot):[1,n] \rightarrow [1,n]$ is a weakly increasing map.
\end{lemma}

\begin{proof}
$\phi(\vec{S} S_{n+1}, n+1) = n+1$ follows from $\phi(\vec{S} S_{n+1}, n+1) \ge n+1$
and $\phi(\vec{S} S_{n+1}, n+1) \in [1,n+1]$. 
By Lemma \ref{lemma-domain-inclusion}.\ref{lemma-domain-inclusion-01}, 
either $\dom(S_i \ldots S_n) = \dom(S_i \ldots S_n S_{n+1})$ 
or $\dom(S_i \ldots S_n) \supset \dom(S_i \ldots S_n S_{n+1})$. 
\begin{enumerate}
\item
\emph{Suppose $\dom(S_i \ldots S_n) = \dom(S_i \ldots S_n S_{n+1})$}.
Then for $j = \phi(\vec{S},i)$
we have $\dom(S_i \ldots S_j) = \dom(S_i \ldots S_n) = \dom(S_i \ldots S_n S_{n+1})$,
while for all $k<j$ we have 
$\dom(S_i \ldots S_k) \not = \dom(S_i \ldots S_n) = \dom(S_i \ldots S_n S_{n+1})$.
This implies $\phi(\vec{S}S_{n+1},i) = \phi(\vec{S},i)$

\item
\emph{Suppose $\dom(S_i \ldots S_n) \supset \dom(S_i \ldots S_n S_{n+1})$}.
Then for all $j \in [i,n]$, $j \ge i$ 
by Lemma \ref{lemma-domain-inclusion}.\ref{lemma-domain-inclusion-02} 
we have 
$\dom(S_i \ldots S_j)  \supseteq \dom(S_i \ldots S_n) \supset \dom(S_i \ldots S_n S_{n+1})$.
We deduce $\phi(\vec{S}S_{n+1},i) >n$ 
and we conclude that $\phi(\vec{S}S_{n+1},i) = n+1$.

\end{enumerate}

Now we check that $\phi$ is weakly increasing.

We assume $i,j \in [1,n]$, $i < j$ in order to prove that $\phi(\vec{S},i) \le \phi(\vec{S},j)$.
Let $b =  \phi(\vec{S},j)$. Then by definition of $\phi$ we have 
$\dom(S_j \ldots S_b) = \dom(S_j \ldots S_n)$. 
By Lemma \ref{lemma-domain-inclusion}.{lemma-domain-inclusion-04} we deduce
$\dom((S_i \ldots S_{j-1}) S_j \ldots S_b) = \dom((S_i \ldots S_{j-1}) S_j \ldots S_n)$.
By definition of $\phi$, this implies $\phi(\vec{S},i) \le b =  \phi(\vec{S},j)$,
as wished.
\end{proof}

 Using the characterization of $\phi(\vec{S},i)$, we have an alternative
way of computing $\minStable(\vec{S})$ by induction on $\vec{S}$.

\begin{lemma}[Computing $\minStable(\vec{S})$ by induction on $\vec{S}$]
\label{lemma-minstable.inductive}
Assume $\vec{S} = S_1, \ldots, S_n$. 
\begin{enumerate}
\item
If $n=0$, then $\minStable(\vec{S}) = \vec{S} = $ the empty sequence.
\item
Suppose $\minStable(\vec{S}) = \vec{T} = T_1, \ldots, T_m$. Then exactly
one of the following holds.
\begin{enumerate}
\item
$\minStable(\vec{S}), S_{n+1}$ is stable and we set
$\minStable(\vec{S}, S_{n+1}) = \minStable(\vec{S}), S_{n+1}$.
\item
$\minStable(\vec{S}), S_{n+1}$ is \emph{not} stable and
$\minStable(\vec{S}, S_{n+1}) = T_1, \ldots, (T_i \ldots T_m S_{n+1})$, 
for the first $i \in [1,m]$ such that 
$\dom(T_i \ldots T_m) \supset \dom(T_i \ldots T_m S_{n+1})$.
\end{enumerate}
\end{enumerate}
\end{lemma}

\begin{proof}
\begin{enumerate}
\item
If $n=0$, then $\minStable(\vec{S}) = \vec{S} = $ the empty sequence by definition
of $\minStable$.

\item
Suppose $\minStable(\vec{S}) = \vec{T} = T_1, \ldots, T_m$.  
By Lemma \ref{lemma-stable-phi-2}
we have $T_i = S_{a_i} \ldots S_{\phi(\vec{S},a_i)}$ with 
$a_1=1$ and $a_{i+1} = \phi(\vec{S},a_i)+1$ 
for all $i \in [1,m-1]$ and $\phi(a_{m}) = n$.  Then exactly one of the following holds.

\begin{enumerate}

\item
Assume $\minStable(\vec{S}), S_{n+1}$ is stable. By 
Lemma \ref{lemma-domain-stability-02}.
\ref{lemma-domain-stability-02},\ref{lemma-domain-stability-04} 
we have
$\dom(T_i \ldots T_m) = \dom(T_i ) = \dom(T_i \ldots T_m S_{n+1})$ for all $i \in [1,m]$,
namely $\dom(S_{a_i} \ldots S_n) = \dom(S_{a_i} \ldots  S_{n+1})$ for all $i \in [1,m]$.
From Lemma \ref{lemma-phi-characterization} we deduce that $\phi(\vec{S},a_i) = 
\phi(\vec{S}S_{n+1},a_i)$ for all $i \in [1,m]$. We have 
$\phi(\vec{S},S_{n+1},\phi(a_{m}) +1)
= \phi(\vec{S},S_{n+1},n+1)=n+1$. By Lemma \ref{lemma-stable-phi-2}, the products of
$\minStable(\vec{S},S_{n+1})$ are $S_{a_i} \ldots S_{\phi(\vec{S},S,a_i)}$, which is
$S_{a_i} \ldots S_{\phi(\vec{S},a_i)}$, namely $T_i$, and the unary product $S_{n+1}$.

We conclude that $\minStable(\vec{S}, S_{n+1}) = T_1, \ldots, T_m, S_{n+1}$.

\item
Assume $\minStable(\vec{S}), S_{n+1}$ is \emph{not} stable. By 
Lemma \ref{lemma-domain-stability-02}.
\ref{lemma-domain-stability-02},\ref{lemma-domain-stability-04} 
we have
$\dom(T_i \ldots T_m) = \dom(T_i ) \not = \dom(T_i \ldots T_m S_{n+1})$ 
for a first $i \in [1,m]$, and by the choice of $i$ that 
$\dom(T_j \ldots T_m) = \dom(T_j ) = \dom(T_j \ldots T_m S_{n+1})$ for all $j \in [1,i-1]$.
For all $j \in [1,i-1]$ we deduce $\dom(T_i \ldots T_m)\supset \dom(T_i \ldots T_m S_{n+1})$.
By Lemma \ref{lemma-domain-inclusion}.\ref{lemma-domain-inclusion-01}
we have $\dom(T_i ) \supseteq  \dom(T_i \ldots T_m S_{n+1})$:
we deduce $\dom(T_i \ldots T_m) \supset \dom(T_i \ldots T_m S_{n+1})$.

By Lemma \ref{lemma-phi-characterization} we have 
$\phi(\vec{S},S_{n+1},a_j) = \phi(\vec{S},a_j)$ for all $j \in [1,i-1]$ and 
$\phi(\vec{S},S_{n+1},a_i) = n+1$. 

By Lemma \ref{lemma-stable-phi-2} and induction on $j \in [1,i-1]$,
the product number $j$ of $\minStable(\vec{S},S_{n+1})$ is $T_j$.
Indeed, assume that $j=1$ or $j>1$ and the previous product is $T_{j-1}$.
In both cases, by Lemma \ref{lemma-stable-phi-2} the first element
of the product number $j$ of $\minStable(\vec{S},S_{n+1})$
is $a_j$, with $a_j = 1$ if $j=1$, and $a_j = \phi(\vec{S},a_ {j-1}+1)$ if $j>1$. 
Again by by Lemma \ref{lemma-stable-phi-2} the product number $j$ 
of $\minStable(\vec{S},S_{n+1})$ is $S_{a_j} \ldots S_{\phi(\vec{S},S_{n+1},a_j)}$, 
which is $S_{a_i} \ldots S_{\phi(\vec{S},a_j)}$, namely $T_j$.
The product number $i$ is the last product and it is $S_{a_i} \ldots S_{n+1}$.

We conclude that 
$\minStable(\vec{S}, S_{n+1}) = T_1, \ldots, T_{i-1}, (T_i \ldots T_m S_{n+1})$.
\end{enumerate}
\end{enumerate}
\end{proof}



\newpage   \section{Injective Domain-Stable Relation Sequences}

We first derive some inequalities for domain-cardinality
valid for sequences of injective relations.

\begin{lemma}[Injective Relation Sequences]
\label{lemma-stable-sequences}
Assume $R, S$, $\vec{R}=R_1, \ldots, R_n$ all are \emph{injective} relations.
\begin{enumerate}
\item
%\label{lemma-stable-sequences-01}
$\Card(\dom(RS)) \le \Card(\dom(S))$
and if $\Card(\dom(RS)) \le \Card(\dom(S))$ then $R$ restricted to $\dom(R) \times \dom(S)$
is a bijection.

\item
%\label{lemma-stable-sequences-02}
For all $i,j \in [1,n]$, $i \le j$:
$$
\Card(\dom(R_i \ldots R_n)) \le \Card(\dom(R_j \ldots R_n))
$$

\item
%\label{lemma-stable-sequences-03}
If  $i, j \in[1,n]$ and $i < j$ and $\Card(\dom(R_i \ldots R_n)) = \Card(\dom(R_j \ldots R_n))$
then $R_i \ldots R_{j-1}$ restricted to $\dom(R_i \ldots R_n) \times \dom(R_j \ldots R_n)$
is a bijection.
\end{enumerate}
\end{lemma}

\begin{proof}
\begin{enumerate}

\item
for all $x \in \dom(RS)$ the sets $X_x = \{ y \in \dom(S) | x R y \}$ are inhabited,
because if $x \in \dom(RS)$ then for some $y \dom(S), z$ we have $x R y z$.
They are pairwise disjoint subsets of $\dom(S)$, because $R$ is injective. We deduce that 
$\Card(\dom(RS)) \le \Card(\bigcup_{x \in \dom(RS)} X_x) \le \Card(\dom(S))$.

Now assume that $\Card(\dom(RS)) = \Card(\dom(S))$. Then 
$\Card(\dom(RS)) = \Card(\bigcup_{x \in \dom(RS)} X_x) = \Card(\dom(S))$. 
From $\Card(\bigcup_{x \in \dom(RS)} X_x) = \Card(\dom(S))$, 
$\bigcup_{x \in \dom(RS)} X_x \subseteq \dom(S)$ and $\dom(S)$ finite we deduce
that $\bigcup_{x \in \dom(RS)} X_x = \dom(S)$. 
From the fact that all $\Card(\dom(RS)) = \Card(\dom(S))$ and all $X_x$ are inhabited we 
deduce that all $X_x$ are singleton. This says that $R$ restricted to $\dom(R) \times \dom(S)$
is a bijection.

\item
Let $i \in [1,n-1]$.
From the point 1 above applied to $R = R_i$, $S = R_{i+1}\ldots R_n$ we deduce that
$\Card(\dom(R_i R_{i+1}\ldots R_n)) \le \Card(\dom(R_{i+1} \ldots R_n))$. The thesis
follows by induction on $j-i$.

\item
$i, j \in[1,n]$ and $i < j$ and $\Card(\dom(R_i \ldots R_n)) = \Card(\dom(R_j \ldots R_n))$.
From the point 1 above applied to $R = R_i \ldots R_{j-1}$, $S = R_i \ldots R_{j-1}$ 
we deduce that $R_i \ldots R_{j-1}$ restricted to 
$\dom(R_i \ldots R_{j-1} \ldots R_j \ldots R_n) \times \dom(R_i \ldots R_{j-1})$
is a bijection.

\end{enumerate}
\end{proof}

We derive more inequalities for domain-cardinality
valid for \emph{stable} sequences of injective relations. 


\begin{lemma}[Injective Domain-Stable Relation Sequences]
\label{lemma-injective-stable-sequences}
Assume $\vec{R}=R_1, \ldots, R_n$ is a \emph{domain-stable}
sequence of \emph{injective} relations.
\begin{enumerate}

\item
\label{lemma-injective-stable-sequences-01}
$$
\forall i,j \in [1,n]. (i \le j) \implies \Card(\dom(R_i)) \le \Card(\dom(R_j))
$$

\item
\label{lemma-injective-stable-sequences-02}
If  $i, j \in[1,n]$ and $i < j$ and $\Card(\dom(R_i)) = \Card(\dom(R_j))$
then $R_i \ldots R_{j-1}$ has domain $\dom(R_i)$ and
when restricted to $\dom(R_i) \times \dom(R_j)$
is a bijection.
\end{enumerate}
\end{lemma}


\begin{proof}
\begin{enumerate}

\item
%\label{lemma-injective-stable-sequences-01}
Assume $i,j \in [1,n]$, $i \le j$. By stability and
Lemma \ref{lemma-domain-stability}.\ref{lemma-domain-stability-04}
we have
$\Card(\dom(R_i)) = \Card(\dom(R_i \ldots R_n))$ and
$\Card(\dom(R_j)) = \Card(\dom(R_j \ldots R_n))$.
By injectivity and Lemma \ref{lemma-stable-sequences}.\ref{lemma-stable-sequences-02}  
we have 
$\Card(\dom(R_i \ldots R_n)) \le \Card(\dom(R_j \ldots R_n))$.
We conclude that 
$\Card(\dom(R_i)) \le \Card(\dom(R_j))$.

\item
%\label{lemma-injective-stable-sequences-02}
Assume $i,j \in [1,n]$, $i < j$. By stability and Lemma \ref{lemma-stable-sequences-03}
$i, j \in[1,n]$ and $i < j$ and $\Card(\dom(R_i \ldots R_n)) = \Card(\dom(R_j \ldots R_n))$
then $R_i \ldots R_{j-1}$ restricted to $\dom(R_i \ldots R_n) \times \dom(R_j \ldots R_n)$.
By stability and
Lemma \ref{lemma-domain-stability}.\ref{lemma-domain-stability-04}
we have
$\dom(R_i) = \dom(R_i \ldots R_n)$ and
$\dom(R_j) = \dom(R_j \ldots R_n)$ and
$\Card(\dom(R_i)) = \Card(\dom(R_i \ldots R_n))$ and
$\Card(\dom(R_j)) = \Card(\dom(R_j \ldots R_n))$.
We conclude that if 
 $i, j \in[1,n]$ and $i < j$ and $\Card(\dom(R_i)) = \Card(\dom(R_j))$
then $R_i \ldots R_{j-1}$ restricted to $\dom(R_i) \times \dom(R_j)$
is a bijection.
\end{enumerate}

\end{proof}


We define a property stronger than domain-stability: constant-cardinality domain.


\begin{definition}[Constant-Cardinality Domain]
Assume $\vec{S}=S_1, \ldots S_n$ is any sequence of relations.
\begin{enumerate}

\item
$\vec{S}$ is a constant-cardinality domain sequence if and only if $\vec{S}$ is is stable
and for all $i,j \in [1,n]$:
$$
\Card(\dom(S_i)) = \Card(\dom(S_j))
$$

\item
If $\vec{S}$ is not empty then we call $\DomCard(\vec{S}) = \Card(S_1)$ 
the domain-cardinality of $\vec{S}$.

\item
A constant-cardinality component $\vec{C}$ of $\vec{S}$ is a maximal segment 
in $\vec{S}$ which is a constant-cardinality domain sequence.
\end{enumerate}
\end{definition}

If $\vec{S}=S_1, \ldots S_n$ is constant-cardinality domain, then $\Card(\dom(S_i)) = \Card(\dom(S_j))$ for all $i,j \in [1,n]$ and for all $k \in [i,n]$, $h \in [j,n]$ by
stability we have
$\Card(\dom(S_i)) =\Card(\dom(S_i \ldots S_k))$ and
$\Card(\dom(S_j)) = \Card(\dom(S_j \ldots S_h))$. This means that all products of
all segments in in $[i,j]$ have the same cardinality: 
$\Card(\dom(S_i \ldots S_k)) = \Card(\dom(S_j \ldots S_h))$.

\begin{lemma}[Decomposing stable sequences of injective relations]
Every stable sequence $\vec{S}$ of injective relations is equal to the 
concatenation $\vec{C_1}, \ldots \vec{C_h}$
of some non-empty constant-cardinality-domain components 
such that $\DomCard(\vec{C_i})  < \DomCard(\vec{C_j})$ for all $i,j \in [1,h]$, $i < j$.
\end{lemma}

\begin{proof}
We define $a_1=a$, $b_i = $ the last $j \in [i,n]$ such that 
$\Card(\dom(S_{a_i})) = \Card(\dom(S_j))$, and if $b_i<n$ then $a_{i+1}=b_i+1$.
Then for some $h \in [1,n]$ we have:

$$
\vec{S}
=
S_{a_1},\ldots,S_{b_1}
\ldots
S_{a_{h}},\ldots,S_{b_{h}}
$$

Each $S_{a_{i}},\ldots,S_{b_{i}}$ is a segment of $\vec{S}$, therefore it is stable
by Lemma \ref{}, and all domains of $S_{a_{i}},\ldots,S_{b_{i}}$ have the same cardinality 
by the choice of $b_i$. Thus, $S_{a_{i}},\ldots,S_{b_{i}}$ is a constant-cardinality domain. 

By the choice of $b_i$ we also have $\Card(S_{b_{i}}) \not = \Card(S_{a_{i+1}})$.
By Lemma \ref{} we deduce that
$\Card(\dom(S_{a_i})) < \Card(\dom(S_{a_j}))$ for all $i,j \in [1,h]$, $i < j$.

\end{proof}

Remark that the subsequence $S_{a_{i}}, \ldots, S_{b_{i}}$ 
are exactly the \emph{maximal} constant-cardinality-domain component of $\vec{S}$.

We describe how the \emph{maximal} constant-cardinality-domain component of $\minStable(\vec{S})$ change as $\vec{S}$ increases.

%17:15 27/01/2026

\begin{lemma}[Constant-cardinality-domain Components and Domain-Stability]
\label{lemma-cardinality-stability}
Assume $\vec{S}, S_{n+1}$ is any sequence of injective relations. Suppose 
the constant-cardinality-domain components of $\minStable(\vec{S})$ and 
$\minStable(\vec{S}, S_{n+1})$ are, respectively, 
$$
\vec{C_1}, \ldots,\vec{C_k} \ \ \ \mbox{ and } \ \ \   \vec{D_1}, \ldots,\vec{D_h}
$$
Then either $h=k+1$ and
$$
\vec{C_j} = \vec{D_j} \ \ \ \mbox{ for all $j \in [1,k]$ }
$$
or for some $i \in [1,h] \cap [1,k]$, and for $c_i = $ the domain-cardinality associated to 
$\vec{C_i}$ and for some colored relation $R$ we have $h=i$ and
$$
\vec{C_j} = \vec{D_j} \ \ \ \mbox{ for all $j \in [1,i-1]$}
$$
and
$$
\mbox{ either } \ \ \ c_i > d_i   \mbox{ and } \vec{D_i} = R
\mbox{ or }      \ \ \ (c_i = d_i) \mbox{ and } \vec{D_i} = \vec{C_i}, R \ \ \ 
$$
\end{lemma}



Lemma \ref{lemma-cardinality-stability} above has the following use. Suppose we have
\begin{enumerate}
\item
a set $\setS$ of sequences of injective relations and $\setC$ the set of  
constant-cardinality-domain components of sequences of $\setS$
\item 
a measure $\mu(\vec{C})$ for each component $\vec{C} \in \setC$, which is 
strictly decreasing w.r.t. $<_1$. 
\end{enumerate}

Then if $\vec{C_1}, \ldots,\vec{C_k}$ are the constant-cardinality-domain components of 
$\minStable(\vec{S})$, with associated cardinalities $c_1, \ldots c_k$,
we take as measure $\mu(\vec{S})$ the list 
$$
c_1,\mu(\vec{C_1}), \ \ \ \ldots, \ \ \ c_k, \mu(\vec{C_k})
$$ 
of measures of components. Lemma \ref{lemma-cardinality-stability} says that that this 
measure is strictly decreasing w.r.t. $<$.


\newpage   \section{Constant-domain of Injective Bicolored Sequences}

If $R= (\Stat,Prog)$ is a bicolored sequence, we define $\universe{R} = \Stat \cup \Prog$
and if $\vec{R}=R_1, ldots, R_n$ is a bicolored sequence, we define
 $\universe{\vec{R}} = \universe{R_1}, \ldots, \universe{R_n}$.

We say that  a bicolored sequence $\vec{R}$ satisfies the injective condition, the domain-stable
condition, the constant domain condition if and only if, respectively, $\universe{\vec{R}}$
satisfies the injective condition, the domain-stable
condition, the constant domain condition. Namely, we move the properties of relation to 
colored relation by forgetting colors.

\begin{definition}[Stationary and progressing relation sequences]
Assume $R$ is an colored relation and $\vec{R}=R_1, \ldots, R_n$ is any colored sequence 
and $i \in [1,n-1]$.
\begin{enumerate}

\item
We say that $R=(\Stat,\Prog)$ is progressing if $\Prog \not = \emptyset$, otherwise we say
that $R$ is 

\item
We say that $\vec{R}$ is progressing in $i$ if and only if 
$R_i \cap (\dom(R_i) \times \dom(R_{i+1}))$ is progressing,
and that that $\vec{R}$ is stationary in $i$ otherwise.
\end{enumerate}
\end{definition}

A constant-domain bicolored sequence $\vec{R}=R_1, ldots, R_n$
of injective relations is stationary if and only if for all $i,j \in [1,n]$, $i < j$ we have
$R_i \ldots R_{j-1} \cap (\dom(R_i) \times \dom(R_j))$ stationary (and necessarily
a stationary bijection).

\begin{lemma}[Constant-domain sequence of injective relations]
\label{lemma-constant-injective}
A stationary constant-domain bicolored sequence $\vec{C}$ of injective bijections 
with associated cardinality $\Card(\dom(C_1))=c$ has length $< \newton{\MaxInd}{c}$.
\end{lemma}


\begin{definition}[Degree of a constant-domain stable colored sequence]
Assume $\setR$ is a set of injective colored relations on $[1,\MaxInd]$ and
$R,S$ is a constant-domain $\setR^+$-sequence. 

We call any $d \in \N$ a degree of the pair $R,S$ 
if and only if for some constant-domain $\setR^+$-sequence 
$R_1, \ldots, R_n, S$ of relations in $\setR^+$ we have:

\begin{enumerate}

\item
$R_1, \ldots, R_n, S$ is a constant-domain $\setR^+$-sequence with
a \emph{stationary} suffix of $d$ elements. 

\item
$R_1 \ldots R_n = R$.

\end{enumerate}
\end{definition}

All pairs have degree $0$, if $R,S$ is stationary in the step $1$ then $d=1$ is a degree
if we choose $n=1$, $R_1 = R$.

A non-empty constant-domain stationary $\setR^+$-sequence $\vec{C}=C_1, \ldots, C_n$, 
with constant cardinality $c$,
has relations with distinct domain $\dom(C_1), \ldots, \dom(C_n)$. There are 
$\newton{\MaxInd}{c}$ subsets of $[1,\MaxInd]$
of cardinality $c$, therefore $n \le \newton{\MaxInd}{c}$. In the same way
we prove that the longest stationary prefix of $\vec{C}$ has at most
$\newton{\MaxInd}{c}$ elements. If $d$ is any of the degrees of
$\vec{C}$, then we can extend $\vec{C}$ in a stationary way to 
an constant-domain stationary stable sequence of injective colored relations for at most
$\newton{\MaxInd}{c}-d$ times. We have a bound to the
number of constant-domain stationary stable injective extensions. $\newton{\MaxInd}{c}-1$
is a bound for all $\vec{C}$, some constant-domain stationary $setR^+$-sequences
have smaller bound, even $0$ when all constant-domain extensions in $\setR^+$
are progressing.

\begin{lemma}[Degree of a constant-domain sequence of injective colored relations]
\label{lemma-degree-constant}
\begin{enumerate}

\item
If $R,S$ has degree $d$,
and $R,S,T$ is constant-domain $\setR^+$-sequence, with $c = \Card(\dom(R))$, then either 
$S \cap \dom(S) \times \dom(T)$ is a progressing bijection on $\dom(S) \times \dom(T)$
or $RS, T$ has degree $d+1$.

\item
If $R,S$ has degree $d$, then $d \le \newton{\MaxInd}{c}$.
\end{enumerate}
\end{lemma}



\newpage   \section{Injective bicolored Relation Sequences}

Assume $\vec{S} = S_1, \ldots, S_n$ is any sequence of bicolored relations. We define
the uncolored version of $\vec{S}$ as $\universe{\vec{S}} = \universe{S_1}, \ldots, 
\universe{S_n}$ and the minimum domain-stable grouping of $\vec{S}$ as
$\minStable(\vec{S}) = \minStable(\universe{\vec{S}})$.

We define the integer sequence $\sigma(i) = \Card(S_i \ldots S_n)$ for $i=1, \ldots, n$.
If $S_1, \ldots, S_n$ are injective,  then $\sigma$ is weakly increasing.
Suppose $\codom(\sigma) = \{c_1, \ldots, c_h\}$ for some increasing sequence 
$0 \le c_1 < c_2 < \ldots < c_h$. For $i=1, \ldots, h$, let $a_i$ be the last index such that
$ \sigma(a_i) = c_i$. If $S_1, \ldots, S_n$ are injective, then $a_i$ is increasing
and $\sigma$ decomposes in $h$ adiacent constant subsequences:
$$
\sigma(1), \ldots, \sigma(a_1),
\ \ \
\sigma(a_1+1), \ldots, \sigma(a_2),
\ \ \ 
\ldots,
\ \ \
\sigma(a_{h-1}+1), \ldots, \sigma(a_h)
$$ 
The subsequence number $i$ has constant value $c_i$.

Suppose:
\begin{enumerate}
\item
$R_1, \ldots, R_n \in \setR^+$ (hence each $R_i$ is injective and satisfies PDC) 
\item
Suppose $R_1, \ldots, R_n$ is stable. 
\end{enumerate}
Set $D_i = \dom(R_i)$ for all $i \in [1,n]$.
%%%%%%%%%%%%%%%%%%%%%%%%%%
% 20:36 24/01/2026
%%%%%%%%%%%%%%%%%%%%%%%%%%
% and $D_n = \codom(R_n)$ if $R_n$ is a bijection 
% and $D_n = \dom(R_n)$ otherwise.
%%%%%%%%%%%%%%%%%%%%%%%%%%
% (i). forse è meglio mettere il caso particolare per D_n tra i possibili miglioramenti
% per la leggibilità della prova ottenuta, anziché nelle definizione di base
%%%%%%%%%%%%%%%%%%%%%%%%%%
% (ii). Altro miglioramento: se il massimo valore del grado \deg(\vec{R},j) per j 
% ultima posizione della componente numero h è M allora si può rimpiazzare 
% 2^\MaxInd e (2^\MaxInd - 1 - \deg(\vec{R},j))  
% con M+1 e con (M - \deg(\vec{R},j)) 
%%%%%%%%%%%%%%%%%%%%%%%%%%%
% (iii). Altro miglioramento: se tutte le immagini della componente di cardinalità h non 
% decrescono nel passaggio nelle componenti successive, si possono eliminare nelle
% componenti successive, perché se decrescono entro la componente di cardinalità h
% allora la componente di cardinalità h decresce e le successive sono irrilevanti.
%%%%%%%%%%%%%%%%%%%%%%%%%%%


\begin{lemma}
\begin{enumerate}
\item
If $i,j \in [1,n]$, $i<j$ and $D_i = D_j$ then $R_i \ldots R_j$ is progressing on $D_i$.
\item
If  $i,j \in [1,n]$, $i<j$ and $\Card(D_i) = \Card(D_j)$ and $j-i \ge 2^{\MaxInd}$ 
then $R_i \ldots R_j$ is progressing on $D_i$.
\end{enumerate}
\end{lemma}

We assume having for each predicate a size map $\size : Assignements, 
[1,\MaxInd] \rightarrow \N$,
such that if $q(\vec{u})$ is the inductive formula number $i$, if $q(\vec{t})$ is the inductive
formula number $j$ and $\sigma$ is an assigment, and $q(\sigma(\vec{u}))$ is an assumption
of the proof of $q(\sigma(\vec{t}))$, then $\size(\sigma,i) < \sigma(\sigma,j)$. We usually 
skip $\sigma$ and we write $\sigma(i), \sigma(j)$.

We include some examples of size map. The first example is the size map counting the
inductive rules in the proof of $q(\vec{t})$.

\begin{enumerate}
\item 
if $q(\vec{t}) = \N(v)$, then $$\sigma(i)=v+1$$ because if $v$ has value $a$ 
then $v=S^a(0)$ and the proof of $\N(v)$ has $a$ times an $S$-rule and $1$-time
the $0$-rule.

\item
If $q(\vec{t}) = \List(l)$ states that $l$ is a list on $\N$, 
then $$\size(i) = (\length(l)+1)+\Sigma_{i=1}^{i=\length(l)}(l[i]+1)$$
Here we count $\length(l)$ times the rule $\cons$, once the rule $\nil$,
the we add the total size of all elements of $l$.

\item
If $q(\vec{t}) = \Tree(t)$ states that $t$ is a binary tree with labels in $\N$, 
then $\size(i)$ is any expression counting the number of rules for nodes in $t$, plus the size
of all labels of $t$.
\end{enumerate}

Alternatively, the size map can count the maximum nesting of inductive rules in the proof of $q(\vec{t})$. In this case the size of $\N(v)$ is $v$ itself.
\\

We claim that after at most $2^\MaxInd$ steps any constant-domain sequence progresses.
We formally express this fact through the notion of degree.

\begin{definition}[Degree of an Index in a Domain-Stable Injective Sequence]
The degree $\deg(j,\vec{R})$ of $j \in [1,n]$ in 
$\vec{R} = R_1, \ldots, R_n$ as above is $2^\MaxInd-1-d$, for $d$ the minimum distance of
an index $i \le j$ such that at least one of the following happens:
\begin{enumerate}
\item
$i=1$
\item
$i>1$ and $\Card(D_{i-1}) < \Card(D_j)$
\item
$R_{i} \ldots R_j$ is progressing between $D_i$ and $D_j$.
\end{enumerate}
If $X \subseteq [1,\MaxInd]$, we set $\size(X) = \Sigma_{i \in X} \size(i)$.
\end{definition}
 
Suppose that in a given node of the proof and for $i \in [1, \MaxInd]$ 
the inductive expression number $i$ is $q(\vec{t})$. Then $\size(i)$ is any expression
of the language computing, in any assignment $\sigma$ to the variable of the node, 
the number of constructors in $\sigma(\vec{t})$. Any map returning a positive integer
and increasing w.r.t. the subterm relation would work. 



If $\vec{R} = R_1, \ldots, R_n \in \setR^+$ is stable then we will define a size for $\vec{R}$: 
$\size(\vec{R})$. We need the stability assumption for $\vec{R}$, 
otherwise we cannot prove that 
the size is a decreasing map. For a generic injective sequence $\vec{S}$, we will first compute
the Minimum Stable Grouping of $\vec{S}$.

\begin{definition}[Size of an Injective Sequence]
If $\dom(R)$ is any injective domain-stable sequence then we set:
$$
\size(\vec{R}) 
= 
(c_1, \size(D_{a_1}), \deg(a_1,\vec{R}), \ldots,c_h,\size(D_{a_h},\deg(a_h,\vec{R}))
$$
If $\dom(S)$ is any injective sequence then we set $\size(\vec{R}) = \size(\minStable(\vec{S}))$.
\end{definition}

If $\vec{R}$ is stable then $\minStable(\vec{R}) = \vec{R}$, therefore 
the general definition includes the definition for the stable case.

The corresponding state for the algorithm translating a cyclic proof into an
injective proof is $\state(\vec{R})$, defined as follows.

\begin{definition}[State of an Injective Domain-Stable Sequence]
If $\dom(R)$ is any injective domain-stable sequence then we set:
$$
\state(\vec{R})
=
(c_1, R_{a_1}\ldots R_{a_2-1}, \deg(a_1,\vec{R}), 
\ldots,
c_h, R_{a_h}\ldots R_{n}, \deg(a_h,\vec{R}))
$$
If $\dom(S)$ is any injective sequence then we set $\state(\vec{R}) = \state(\minStable(\vec{S}))$.
\end{definition}
 
We prove that adding any $R_{n+1}$ to the sequence reduces the index in the 
lexicographic ordering: $\size(\vec{R}) > \size(\vec{R}R_{n+1}) $. Besides,
$\size(\vec{R})$ can be computed from $\state(\vec{R})$. Indeed, if
$\vec{R} = \minStable(\vec{S})$, then $\vec{R}$ is stable and

$$
\size(\vec{S})
=
\size(\vec{R})
=
(
 c_1,(\dom(R_{a_1}\ldots R_{a_2-1}), \deg(a_1,\vec{R}), 
 \ldots,
 c_h,\dom(R_{a_h}\ldots R_{n}), \deg(a_h,\vec{R})
 )
$$

We can compute $\state(\vec{S}S_{n+1})$ from $\state(\vec{S})$.
There are six cases. Suppose $\minStable(\vec{S}) = \vec{R}$.
We have $\Card(\dom(S_{n+1}) \ge c_h$.

\begin{enumerate}
\item
Suppose$ \vec{R}S_{n+1}$ is stable and $\Card(\dom(R_{n+1})> c_h$. Then
$c_{h+1} = \Card(\dom(R_{n+1})$, $a_h = n$, $a_{h+1} = n+1$ and 
$\state(\vec{S}S_{n+1})$ is 

$$
(c_1, R_{a_1}\ldots R_{a_2-1}, \deg(a_1,\vec{R}), 
\ldots,
c_h, R_{a_h}, \deg(a_h,\vec{R}),
c_{h+1}, S_{n+1}, 2^{\MaxInd} - 1)
$$

\item
Suppose$ \vec{R}S_{n+1}$ is stable and $\Card(\dom(R_{n+1})= c_h$
and $R_n S_{n+1}$ is progressing on $\dom(R_{n})$.
Then the new value of $a_h$ is $n+1$ and $\state(\vec{S}S_{n+1})$ is 

$$
(c_1, R_{a_1}\ldots R_{a_2-1}, \deg(a_1,\vec{R}), 
\ldots,
c_{h-1}, R_{a_{h-1}}\ldots R_{a_h}, \deg(a_h,\vec{R}),
c_{h}, S_{n+1}, 2^{\MaxInd} - 1)
$$

\item
Suppose$ \vec{R}S_{n+1}$ is stable and $\Card(\dom(R_{n+1})= c_h$
and $R_n S_{n+1}$ is stationary on $\dom(R_{n})$.

Then $\state(\vec{S}S_{n+1})$ is 

$$
(c_1, R_{a_1}\ldots R_{a_2-1}, \deg(a_1,\vec{R}), 
\ldots,
c_{h-1}, R_{a_{h-1}}\ldots R_{a_h}, \deg(a_h,\vec{R}),
c_{h}, S_{n+1}, \deg(a_h,\vec{R}) - 1)
$$

\item
Suppose$ \vec{R}S_{n+1}$ is not stable and $j \in [1,n]$ is the first index
such that $\dom(R_j \ldots R_n) \supset \dom(R_j \ldots R_n S_{n+1})$.
Then $j = a_i + 1$ for some $i$ and the Minimum Stable grouping of $\vec{S}S_{n+1}$
is $R_1, \ldots, R_{a_i}, R$ with $R = R_j \ldots R_n S_{n+1}$. If $R_{a_i} R$
is progressing on $\dom(R_{a_i})$ then the new value of $a_i$ is
$a_{i+1}$ and $\state(\vec{S}S_{n+1})$ is 

$$
(c_1, R_{a_1}\ldots R_{a_2-1}, \deg(a_1,\vec{R}), 
\ldots,
c_{i-1}, R_{a_{i-1}}\ldots R_{a_i}, \deg(a_i,\vec{R}),
c_{i}, R, 2^{\MaxInd}-1)
$$

\item
If instead $R_{a_i} R$
is stationary on $\dom(R_{a_i})$ then $\state(\vec{S}S_{n+1})$ is:

$$
(c_1, R_{a_1}\ldots R_{a_2-1}, \deg(a_1,\vec{R}), 
\ldots,
c_{i-1}, R_{a_{i-1}}\ldots R_{a_i}R, \deg(a_i,\vec{R})-1)
$$

\end{enumerate}


The number of states is bounded by 
$
(2^{O(\MaxInd^2)} \times 2^\MaxInd \times \MaxInd)
\times \ldots \times
(2^{O(\MaxInd^2)} \times 2^\MaxInd \times \MaxInd)
$,
with at most $3\MaxInd$ product. 
We conclude that it is bounded by $2^{O(\MaxInd^3)}$. Each step requires time bounded by
$O(2^{\MaxInd^2}) \times \log_2(2^{O(\MaxInd^3)})$. 
We have $n \le 2^{\MaxInd^2}$, and cyclically 
for each $R_i$ with $i = 1, \ldots, n$ we add the state obtained by concatenating with 
$R_i$ and we compare it with the existing states sorted, in order to decide whether it is a new
state or not. We stop adding states in the first cycle we add no new state.
The total time cost is bounded by the number of relations times the cost of checking
whether we have a new state times the upper bound to the number of states.
The total is $O(2^{\MaxInd^2}) \times {O(\MaxInd^3)} \times 2^{O(\MaxInd^3)}$,
namely it is $2^{O(\MaxInd^3)}$ again. Therefore the total cost is poly-exponential.

%If $\vec{R} = R_1, \ldots, R_n \in \setR^+$ and $\vec{S}$ is the Minimum Stable
%grouping of $\vec{R} $ then we set $\size(\vec{R})  = \size(\vec{S}) $.

\newpage   \section{A Ranking Function}
We define a ranking function. Assume $\vec{S}$ has minimum stable grouping 
$\minStable{\vec{C}} = \vec{R} = R_1, \ldots, R_n$. Assume there are $h$ maximal
constant-cardinality domain components, and the component number $i$ is  
$R_{a_i}, \ldots ,R_{b_i}$, with $c_i  = \Card(R_{a_i})$.  

We associate to each colored $R$ and each node $\alpha$ an index $\size(\alpha, R) =
\Sigma_{x \in \dom(R)} \max(\size(R(x))$, with $\size(X)=\Sigma_{x \in X} \size(x)$
and if the inductive formula number $x$ of $\alpha$ is $\N(t)$, then $\size(x) =t$.

\begin{lemma}[Size of a Colored Relation Eventually Decreases]
\label{lemma-size-colored-relation}
Assume $\setR^+$ satisfies PDC and the Injectivity Condition IC. 
Assume $R_1, \ldots, R_n, R_{n+1} \setR$ is a constant-cardinality-domain sequence of 
relations in $\setR^+$. Then:

$$
\size(R_1 \ldots R_n) = \size(R_1 \ldots R_n R_{n+1})
\implies
R_n \cap \dom(R_{n}) \times \dom(R_{n+1}) \mbox{ stationary }
$$
\end{lemma}

\begin{proof}
\ldots\ldots\ldots
\end{proof}

\begin{corollary}[Length Bound from PDC, IC, Constant-Cardinality Domain]
\begin{enumerate}
\item
For all colored relations $R$: $size(\dom(R)) \ge \size(R)$.
 
\item
If $R_1 \cap (\dom(R_1) \times \dom(R_2))$ then $\size(R_1) \ge \size(\dom(R_2))$.

\item
Assume $R_1,R_2$ is a constant-cardinality domain sequence 
of injective colored relations.
If $\size(R_1) = \size(R_2)$, then $\size(R_1) = \size(\dom(R_2)) = \size(R_2)$ and 
$R_1 \cap (\dom(R_1) \times \dom(R_2))$ is a
stationary bijection.

\item
Assume $R_1, \ldots, R_n, R_{n+1}$ is a constant-cardinality-domain sequence of 
relations in $\setR^+$ with domain-cardinality $c$.
Suppose for some $k$
we have $\size(R_i) = k$ for all $i \in [1,n+1]$. Then $\size(\dom(R_i)) = k$
for all $i \in [2,n+1]$ and 
$R_i \cap (\dom(R_i) \times \dom(R_{i+1}))$ is a stationary bijection for all $i \in [2,n]$ and
$$
n \le \newton{\MaxInd}{c}+1
$$ 
\end{enumerate}
\end{corollary}


\begin{proof}

\begin{enumerate}
\item
For all $x \in \dom(R)$, $y \in R(x)$ from $x R y$ and the conditions 
about coloring, this implies that $\size(x) \ge size(y)$. Thus, $\size(x) \ge \max(\size(R(x))$.
We conclude that $size(\dom(R)) = \Sigma_{x \in \dom(R)} \size(x) \ge 
\Sigma_{x \in \dom(R)} \max(\size(R(x)) = \size(R)$.

\item
If $R_1 \cap (\dom(R_1) \times \dom(R_2))$ is a bijection.
then for all $y \in \dom(R_2)$ there is exactly one $x \in \dom(R_1)$
such that $y \in R_1(x)$. The conditions about coloring imply that $\size(x) \ge size(y)$. 
Thus, $ \size(R_1) = \Sigma_{x \in \dom(R_1)} \max(\size(R_1(x)) \ge 
\Sigma_{y \in \dom(R_2)} \size(y) = \size(\dom(R_2)$.

\item
By $R_1, R_2$ constant-cardinality-domain sequence and $R_1, R_2$ injective and
Lemma \ref{} we deduce that $R_1 \cap (\dom(R_1) \times \dom(R_2))$ is a bijection.
By points above we deduce that $\size(R_1) \ge \size(\dom(R_2)) \ge \size(R_2)$,
and from $\size(R_1) = \size(R_2)$ we deduce 
$\size(R_1) = \size(\dom(R_2)) = \size(R_2)$. 

\item
By Lemma \ref{lemma-size-colored-relation} $R_i \cap \dom(R_i) \times \dom(R_{i+1})$ 
is a stationary bijection for all $i \in [1,n ]$. By the previous point we have 
$\size(\dom(R_i)) = k$ for all $i \in [2,n+1]$. Thus, $R_i \cap (\dom(R_i) \times 
\dom(R_{i+1}))$ is a stationary bijection for all $i \in [2,n]$, otherwise we would have
$\size(\dom(R_i)) > \size(\dom(R_{i+1}))$
By Lemma \ref{} implies that the sets $\dom(R_i)$ are distinct 
for $i \in [2,n ]$, namely that they are at most $\newton{\MaxInd}{c}$.
We conclude that $(n-2)+1 \le  \newton{\MaxInd}{c}$, that is, 
$n \le \newton{\MaxInd}{c}+1$.
\end{enumerate}
\end{proof}


We define $Eq(\alpha, \beta)$ the set of equations $\size(x) > \size(y)$ such that the
the inductive formula number $x$ of $\alpha$ progresses to the inductive formula number $y$
 of $\beta$.
 
 We associate to $\alpha$ and $\vec{S}$ a list 
 $f(\alpha,\vec{S})=(l_1, \ldots, l_h)$ elements, with the element 
 number $i$ equal to 
 $(
 c_i, 
 \size(\alpha, R_{b_i}\ldots, R_n), 
 \newton{\MaxInd}{c_i} - \deg(R_{a_i} \ldots R_{b_i-1}, R_{b_i})
 )$.
 

We claim that if $\beta$ is a bud with companion $\alpha$, then 
$Eq(\alpha, \beta) \vdash f(\alpha,\vec{S}) > f(\beta,\vec{S})$

The minimum of all $f(\alpha,\vec{S})$ for all $\vec{S}$ is the required ranking function.
We represent $\vec{S}$ with $h$ pairs of the form $(R_{a_i}\ldots R_{b_i-1}),R_{b_i}$




 


\end{document}

